\documentclass[11pt,a4paper]{article}
\usepackage[spanish,es-tabla]{babel}
\usepackage[margin=2.4cm]{geometry}
\usepackage{amsmath,amssymb,amsthm,bm,booktabs,graphicx,xcolor,mathtools}
\newtheorem{fact}{Fact}
\newtheorem{worked}{Worked example}
\newcommand{\Cof}{\operatorname{cof}}
\newcommand{\tr}{\operatorname{tr}}
\newcommand{\R}{\mathbb{R}}
\newcommand{\SO}{\mathrm{SO}(2)}

\newenvironment{takeaway}
 {\begin{quote}\small\itshape\color{blue!60!black}}
 {\end{quote}}

\title{Polyconvexity, los invariants, y por qu\'e construimos los nuestros as\'i\\
\large Versi\'on Spanglish del documento did\'actico sobre la construcci\'on anisotropic del PANN}
\author{}
\date{\today}

\begin{document}
\maketitle

\begin{abstract}
\textbf{Nota del traductor.} Este documento mantiene en \textbf{ingl\'es}
los t\'erminos t\'ecnicos centrales tal como aparecen en la literatura
(\emph{polyconvex/polyconvexity}, \emph{convex/convexity},
\emph{objectivity}, \emph{invariant(s)}, \emph{ICNN}, \emph{Green--Lagrange
strain}, \emph{deformation gradient}); el resto est\'a en espa\~nol
natural. Si prefieres que algo espec\'ifico tambi\'en quede en ingl\'es
(o al rev\'es), dime y lo ajusto.

Este es un documento did\'actico, no un paper. Su \'unico objetivo es que
entiendas, \emph{de verdad} y desde cero, qu\'e es \emph{polyconvexity},
por qu\'e la hyperelasticity necesita esta condici\'on tan espec\'ifica
(y un poco rara a primera vista) en lugar de la \emph{convexity} normal,
qu\'e son en realidad los ``tres \emph{invariants} famosos'' $I_1,I_2,I_3$
y por qu\'e son \emph{polyconvex}, y --la parte que es genuinamente
espec\'ifica de este proyecto-- qu\'e son nuestras propias $15$
features, por qu\'e tienen esa forma, si son una generalizaci\'on de
$I_1,I_2,I_3$ o algo completamente distinto, y de d\'onde sali\'o la
idea. Cada ejemplo num\'erico de abajo fue calculado y verificado, no
inventado, as\'i que puedes reproducir cada n\'umero t\'u mismo.
\end{abstract}

\tableofcontents

\section{¿Por qu\'e nos importa la convexity, para empezar?}
\label{sec:motivation}

Empecemos por la pregunta m\'as b\'asica posible, sin nada de
elasticity todav\'ia: ¿por qu\'e la \emph{convexity} es tan importante
en optimizaci\'on y en f\'isica?

\begin{takeaway}
Una funci\'on \emph{convex} no tiene m\'inimos locales ``falsos'' en los
que quedarte atascado, y una l\'inea que conecta dos puntos cualquiera
de su gr\'afica siempre queda por encima de la gr\'afica. Este \'unico
hecho geom\'etrico es lo que hace que los problemas de minimizaci\'on
construidos a partir de funciones \emph{convex} se comporten bien:
los solvers basados en gradiente (incluyendo el m\'etodo de Newton)
convergen de forma confiable, y la existencia de un minimizador est\'a
garantizada bajo condiciones adicionales bastante suaves.
\end{takeaway}

Toma el ejemplo m\'as simple posible, $f(x)=x^2$. Tiene un solo m\'inimo,
en $x=0$, y el m\'etodo de Newton lo encuentra en un solo paso desde
cualquier punto de partida. Ahora compara $g(x)=x^4-3x^2$ (una forma de
``doble pozo''): tiene \emph{dos} m\'inimos locales y un m\'aximo local
en medio. Un solver basado en gradiente que arranque cerca del pozo
``equivocado'' converger\'a tranquilamente a un minimizador que no es
el global, y si est\'as minimizando una energ\'ia que se supone
representa un equilibrio f\'isico real, converger al m\'inimo local
equivocado es una respuesta f\'isicamente incorrecta de verdad, no
simplemente un inconveniente num\'erico.

En hyperelasticity, el ``problema de optimizaci\'on'' es: encontrar el
campo de desplazamientos que minimiza la energ\'ia el\'astica total,
sujeto a las condiciones de frontera. Si la densidad de energ\'ia $W$
(como funci\'on de la deformaci\'on) es \emph{convex}, resultados
cl\'asicos (el \emph{direct method of the calculus of variations}) te dan
la existencia de un minimizador casi gratis, y el m\'etodo de Newton
sobre el problema discretizado tiende a comportarse bien. Esta es toda
la motivaci\'on para querer alg\'un tipo de \emph{convexity}. La
pregunta es: \emph{convexity} ¿de qu\'e, exactamente?

\section{La respuesta ingenua falla: $W$ no puede ser \emph{convex} en $F$}
\label{sec:naive-fails}

Lo m\'as ingenuo que se puede pedir es que $W$, como funci\'on del
\emph{deformation gradient} $F$ (una matriz $2\times2$, en nuestro caso
2D), sea \emph{convex} en $F$. Esto falla, y falla por una raz\'on
limpia y verificable que no tiene nada que ver con ning\'un material en
particular: la \emph{objectivity} (indiferencia al marco de referencia)
lo hace imposible.

La \emph{objectivity} dice que $W(QF)=W(F)$ para toda rotaci\'on
$Q\in\SO$: rotar r\'igidamente el cuerpo deformado, sin deformarlo m\'as,
no puede cambiar la energ\'ia almacenada. Toma $Q=-I$ (una rotaci\'on de
$180^\circ$, que es un elemento genuino de $\SO$ en dos dimensiones). La
\emph{objectivity} obliga a que
\[
 W(-F)=W(F).
\]
Ahora bien, si $W$ fuera \emph{convex} en $F$, el valor en el punto
medio del segmento entre $F$ y $-F$ tendr\'ia que ser, como m\'aximo,
el promedio de los extremos:
\[
 W\!\left(\tfrac12F+\tfrac12(-F)\right) = W(\bm 0) \;\leq\; \tfrac12W(F)+\tfrac12W(-F) = W(F).
\]
Pero el punto medio es la \emph{matriz cero}, $F=\bm 0$, que tiene
$\det F=0$: f\'isicamente, este es un estado en el que el material ha
sido aplastado hasta volumen cero, y cualquier energ\'ia f\'isicamente
sensata tiene que explotar ah\'i ($W\to\infty$), no quedarse finita y
por debajo de $W(F)$. La \emph{convexity} en $F$ y la \emph{objectivity}
son sencillamente incompatibles. Esto no es un defecto de ning\'un
modelo en particular; es un teorema (se remonta a Coleman y Noll, y se
discute en pr\'acticamente todo tratamiento moderno del tema, por
ejemplo el libro de texto de Ciarlet~\cite{Ciarlet1988}).

\begin{takeaway}
No puedes pedir que $W$ sea una funci\'on \emph{convex} ordinaria de
$F$. Cualquier noci\'on de \emph{convexity} utilizable para elasticity
tiene que ser algo m\'as d\'ebil, compatible con $W(-F)=W(F)$ y con
$W\to\infty$ cuando el material es aplastado.
\end{takeaway}

\section{El truco de Ball: levantar el problema a un espacio m\'as grande}
\label{sec:lifting-trick}

Aqu\'i est\'a la idea central, y vale la pena entenderla una vez, en
abstracto, completamente separada de la elasticity, porque despu\'es la
versi\'on de elasticity deja de parecer misteriosa.

\begin{takeaway}
Una funci\'on puede no ser \emph{convex} en su variable
\emph{original}, y sin embargo volverse \emph{convex} en el momento en
que le das, como si fuera un input independiente, alguna cantidad
\emph{extra, derivada} de la original.
\end{takeaway}

\begin{worked}[El truco del ``lifting'', sin nada de elasticity]
\label{ex:lifting}
Sea $F=\begin{psmallmatrix}a&b\\c&d\end{psmallmatrix}$ y considera la
funci\'on $h(F)=\det F=ad-bc$. Como funci\'on de los cuatro n\'umeros
$(a,b,c,d)$, $h$ es una forma bilineal con forma de silla de montar: su
Hessian es indefinido (tiene autovalores positivos y negativos), as\'i
que $h$ \emph{no} es \emph{convex} en $F$. Mover $b$ y $c$ en la
direcci\'on correcta puede hacer que $h$ baje por debajo del promedio de
dos puntos donde es grande --exactamente lo que la \emph{convexity}
proh\'ibe.

Ahora hagamos un truco: introduzcamos un s\'imbolo nuevo, extra,
$\delta$, declaremos $\delta:=\det F$, y consideremos la funci\'on del
\emph{par} $(F,\delta)$ dada por
\[
 \tilde h(F,\delta) := \delta.
\]
¡Esto ahora es trivialmente \emph{convex} (de hecho es af\'in/lineal) en
el par $(F,\delta)$, porque ya no depende de $F$ en absoluto una vez
que $\delta$ se trata como una coordenada independiente! Por supuesto,
$\tilde h(F,\det F) = h(F)$ recupera la funci\'on original, no
\emph{convex} --nada f\'isico cambi\'o. Lo que cambi\'o es que ahora
podemos hacer una pregunta \emph{distinta, m\'as f\'acil}: ¿es $\tilde
h$ conjuntamente \emph{convex} en el par ampliado de variables
independientes $(F,\delta)$, en lugar de en $F$ solo, restringido al
conjunto curvo $\{(F,\det F)\}$? La respuesta puede ser s\'i incluso
cuando la respuesta a la pregunta original es no.
\end{worked}

Esto es, casi literalmente, lo que hace Ball~\cite{Ball1977} para la
elasticity. En vez de preguntar si $W$ es \emph{convex} en $F$,
pregunta si existe una funci\'on $G$ de \emph{tres variables de bloque
independientes} $(F,\Cof F,\det F)$ --tratando la matriz cofactor
$\Cof F$ y el determinante $\det F$ como si fueran inputs extra,
independientes, exactamente como $\delta$ arriba-- tal que
\begin{equation}
 W(F)=G(F,\Cof F,\det F),\qquad G\text{ conjuntamente \emph{convex} en }
 \R^{2\times2}\times\R^{2\times2}\times\R.
 \label{eq:polyconvex-def}
\end{equation}
Esto es \emph{polyconvexity}. Es estrictamente m\'as d\'ebil que la
\emph{convexity} en $F$ (toda funci\'on \emph{convex} en $F$ es
trivialmente \emph{polyconvex}, ignorando los bloques extra, pero no al
rev\'es), es compatible con la \emph{objectivity} (porque
$\Cof(-F)=\Cof F$ y $\det(-F)=\det F$ en 2D, as\'i que
$G(-F,\Cof F,\det F)$ no tiene por qu\'e ser igual a
$G(F,\Cof F,\det F)$ por la misma raz\'on por la que a $\delta$ no le
importaba el truco del cambio de signo de arriba --la parte no
\emph{convex}, sensible al signo, del problema es exactamente la parte
que qued\'o ``absorbida'' en los bloques extra), y es exactamente la
condici\'on que Ball necesitaba, junto con una \emph{growth condition},
para probar la existencia de minimizadores de energ\'ia mediante el
\emph{direct method}. No dice nada, por s\'i sola, sobre el signo de
$\partial^2W/\partial F^2$ en una deformaci\'on dada --esa es una
pregunta completamente distinta, y estrictamente m\'as fuerte, que se
discute en el paper principal.

\section{Los tres \emph{invariants} famosos}
\label{sec:classical-invariants}

Vamos directo a tu pregunta: s\'i, has o\'ido hablar de tres
\emph{invariants} famosos, y aqu\'i est\'an. Para el right
Cauchy--Green tensor $C=F^TF$ (sim\'etrico, definido positivo), los
\emph{invariants} cl\'asicos usados en pr\'acticamente todo libro de
texto de hyperelasticity (ver Ogden~\cite{Ogden1984}, o
Ciarlet~\cite{Ciarlet1988}) son
\begin{align}
 I_1 &= \tr C, \label{eq:I1}\\
 I_2 &= \tfrac12\big[(\tr C)^2-\tr(C^2)\big]
       = \tr\!\big(\Cof(C)\big)\quad\text{(en 3D)}, \label{eq:I2}\\
 I_3 &= \det C = J^2. \label{eq:I3}
\end{align}
F\'isicamente: $I_1$ es (el doble de) la suma de los cuadrados de los
principal stretches, as\'i que mide cu\'anto se ha estirado el material
en total, combinando todas las direcciones; $I_3=J^2$ mide el cambio de
volumen al cuadrado; $I_2$ tiene que ver con cu\'anto han cambiado las
\emph{\'areas} (en lugar de las longitudes), y por eso se construye a
partir de $\Cof(C)$, el objeto que describe c\'omo se transforman los
elementos de \'area orientados bajo $F$.

Los dos modelos hyperelastic compresibles/incompresibles m\'as famosos
que existen se construyen directamente a partir de estos:
\[
 W_{\rm neo\text{-}Hookean}=c_1(I_1-3), \qquad
 W_{\rm Mooney\text{-}Rivlin}=c_1(I_1-3)+c_2(I_2-3),
\]
y ambos son \emph{polyconvex}, por una raz\'on estructural muy limpia
que ya puedes ver: $I_1=\tr C=|F|^2$ (una suma de cuadrados de las
entradas de $F$) es \emph{convex} en $F$ directamente --vive
enteramente en el \emph{primer} bloque de~\eqref{eq:polyconvex-def}-- e
$I_2$, al estar construido a partir de $\Cof(C)$, vive en el
\emph{segundo} bloque. Para coeficientes no negativos $c_1,c_2\geq0$,
una combinaci\'on no negativa de funciones que son cada una \emph{convex}
en su propio bloque vuelve a ser conjuntamente \emph{convex} en la
variable combinada $(F,\Cof F,J)$ --exactamente la regla ``suma de
piezas \emph{convex} es \emph{convex}'' usada en todo el paper
principal.

\begin{fact}[Una sutileza de 2D que vale la pena conocer]
\label{fact:2d-degeneracy}
En tres dimensiones, $I_1,I_2,I_3$ son tres n\'umeros genuinamente
independientes. En dos dimensiones --nuestro caso-- $C$ solo tiene dos
autovalores $\lambda_1,\lambda_2$ (principal stretches al cuadrado), y
el rol de ``$I_2$'' (suma de productos de pares de autovalores) colapsa
sobre un \'unico par, $\lambda_1\lambda_2=\det C=I_3$. Entonces en 2D,
el $I_2$ cl\'asico y el $I_3$ cl\'asico son literalmente el mismo
n\'umero, y la imagen ingenua de ``tres \emph{invariants}'' en realidad
solo tiene dos escalares independientes, $I_1$ y $J$. Esto es
exactamente por qu\'e la teor\'ia 2D de Ball se formula usando toda la
\emph{matriz} $\Cof F$ como bloque, no solo su traza escalar: tratar
$\Cof F$ como un objeto $2\times2$ completo (no meramente ``el segundo
\emph{invariant}'') deja espacio para extraer informaci\'on
genuinamente nueva, no isotropic, de \'el --cosa que su sola traza
nunca podr\'ia hacer (su traza, como acabamos de ver, solo te devuelve
$I_1$). Este es precisamente el espacio que usa nuestra propia
construcci\'on m\'as abajo.
\end{fact}

\section{¿Usamos $I_1,I_2,I_3$? -- S\'i, generalizados}
\label{sec:are-we-using-them}

Aqu\'i est\'a la respuesta directa a ``¿usamos los tres \emph{invariants}
famosos?'', hecha precisa en lugar de vaga.

\begin{fact}[$I_1$ es un caso muy particular de nuestra propia construcci\'on]
\label{fact:i1-special-case}
Una identidad b\'asica y exacta: para \emph{cualquier} base ortonormal
$\{e_1,e_2\}$ de $\R^2$ (no solo la est\'andar),
\[
 I_1=\tr C = e_1^TCe_1+e_2^TCe_2 .
\]
Es decir, $I_1$ ya es, en secreto, una suma pesada de $d\cdot Cd$ sobre
dos direcciones unitarias elegidas, con peso $1$ en cada una y potencia
$p=1$. Compara esto con nuestra propia familia de features (paper
principal, ecuaci\'on de $Q^F_{k,p}$),
\[
 Q_{k,p}^F=\sum_{i=1}^3 w_{ki}\,(d_{ki}\cdot Cd_{ki})^p .
\]
$I_1$ es exactamente el caso particular de esta f\'ormula con: solo
\emph{dos} direcciones en lugar de tres, esas dos direcciones
\emph{ortogonales} entre s\'i, pesos \emph{iguales} a $1$, y potencia
$p=1$. Nuestra construcci\'on no es an\'aloga a $I_1$; es una
generalizaci\'on literal y directa de $I_1$ --m\'as direcciones (tres,
no dos), sin requerir que sean ortogonales, pesos positivos generales
(sujetos a la \emph{balance condition} de m\'as abajo, que es
exactamente lo que reemplaza a ``ortonormal'' una vez que renuncias a
la ortogonalidad), y potencias m\'as altas $p=2,3$ en lugar de $p=1$.
\end{fact}

La misma historia se repite un nivel arriba: nuestras features
$Q_{k,p}^H$ usan $d\cdot\Cof(C)d$ exactamente de la misma forma en que
$Q^F$ usa $d\cdot Cd$, as\'i que generalizan el \emph{invariant} basado
en $\Cof$ (con sabor a $I_2$) exactamente en el mismo sentido. Y $J,J^2$
no se generalizan en absoluto --son $I_3$ e $I_3^2$, usados sin ning\'un
cambio, porque el \emph{invariant} volum\'etrico no necesita una
generalizaci\'on anisotropic: el cambio de volumen ya es un \'unico
escalar independiente de la direcci\'on, y no hay nada anisotropic que
generalizar ah\'i.

\begin{takeaway}
Entonces: s\'i, estamos usando $I_1$, $I_2$ (v\'ia $\Cof$), e $I_3=J^2$
--pero no como tres n\'umeros pelados. Usamos \emph{muchas} versiones
direccionales de las construcciones tipo-$I_1$ y tipo-$I_2$ (tres
direction systems, cada uno contribuyendo t\'erminos en potencias
$p=2,3$), porque un material \emph{sin ninguna symmetry asumida}
necesita poder responder de forma distinta en direcciones distintas, y
un solo escalar $I_1$ o $I_2$, al estar construido por suma sobre una
base ortonormal (o, de hecho, \emph{cualquier} base), es completamente
ciego a la direcci\'on por construcci\'on.
\end{takeaway}

\section{Por qu\'e $I_1,I_2,I_3$ solos no pueden describir un material anisotropic}
\label{sec:isotropy-blindness}

Esto hay que verlo num\'ericamente para que realmente se entienda, as\'i
que aqu\'i est\'a, calculado de verdad, no simplemente afirmado.

\begin{worked}[Dos materiales distintos que $I_1$ y $J$ no pueden distinguir]
\label{ex:isotropy-blind}
Toma dos estados de deformaci\'on,
\[
 F_a=\begin{psmallmatrix}1.4&0\\0&0.9\end{psmallmatrix}
 \quad(\text{estirar }1.4\times\text{ en }x,\ \text{comprimir a }0.9\times\text{ en }y),
 \qquad
 F_b=\begin{psmallmatrix}0.9&0\\0&1.4\end{psmallmatrix}
 \quad(\text{roles de }x,y\text{ intercambiados}).
\]
Estos son estados f\'isicamente muy distintos para un material
anisotropic --uno estira fuerte en $x$ y comprime en $y$, el otro hace
lo contrario-- y un material cuya rigidez depende de la direcci\'on
deber\'ia, en general, almacenar una cantidad distinta de energ\'ia en
cada caso. Sin embargo:
\[
 I_1(F_a)=I_1(F_b)=2.77, \qquad J(F_a)=J(F_b)=1.26 .
\]
Exactamente iguales, hasta el \'ultimo decimal, porque la traza y el
determinante literalmente no saben cu\'al eje es cu\'al. \emph{Cualquier}
energ\'ia construida puramente a partir de $I_1,I_2,I_3$ (es decir,
cualquier modelo hyperelastic \emph{isotropic}, sin importar c\'omo se
combinen los \emph{invariants}) es matem\'aticamente incapaz de
distinguir $F_a$ de $F_b$. Esto no es un error de aproximaci\'on; es
exacto y estructural.
\end{worked}

\section{Nuestra construcci\'on: direction systems, \emph{balance}, y por qu\'e la potencia c\'ubica no es decorativa}
\label{sec:our-construction}

Para arreglar el Ejemplo~\ref{ex:isotropy-blind}, necesitas al menos una
feature que \emph{no} sea ciega a la direcci\'on. Para eso son los tres
direction systems (tabla de direction systems en el paper principal).
Cada sistema $k$ consiste en tres vectores unitarios $d_{k1},d_{k2},
d_{k3}$ en \'angulos espec\'ificos, con pesos positivos espec\'ificos
$w_{k1},w_{k2},w_{k3}$, que satisfacen la \emph{balance condition}
\begin{equation}
 \sum_{i=1}^3 w_{ki}\,d_{ki}\otimes d_{ki} = I .
 \label{eq:balance}
\end{equation}
Esta condici\'on es la generalizaci\'on directa en 2D de ``ortonormal
con pesos unitarios'': dice que las direcciones elegidas (posiblemente
obl\'icuas, no ortogonales), pesadas apropiadamente, siguen sumando la
identidad, de la misma forma en que $e_1\otimes e_1+e_2\otimes e_2=I$
lo hace para un par ortonormal. Esto es exactamente lo que hace que la
prueba de \emph{stress-free reference state} (C4 en el paper principal)
funcione con una sola correcci\'on escalar, y es por qu\'e el modelo se
comporta \emph{isotropically en el estado no deformado} aunque sea
completamente anisotropic lejos de \'el --un material real con una
textura o un tejido normalmente se ve localmente isotropic justo en
reposo, y solo revela su anisotropy una vez que lo deformas, que es
exactamente el comportamiento que esta construcci\'on reproduce por
dise\~no.

\begin{worked}[Calculando $Q^F_{1,2}$ y $Q^F_{1,3}$ a mano, de verdad]
\label{ex:worked-full}
El sistema~1 usa \'angulos $0^\circ,50^\circ,130^\circ$ con pesos
$0.2959,0.8520,0.8520$. Para $F_a$ de arriba, $C=\mathrm{diag}(1.96,0.81)$.
Las tres direcciones son $d_1=(1,0)$, $d_2=(0.6428,0.7660)$,
$d_3=(-0.6428,0.7660)$, lo que da
\[
 d_1\cdot Cd_1=1.96,\qquad d_2\cdot Cd_2=d_3\cdot Cd_3=1.2852 .
\]
Entonces
\[
 Q^F_{1,2}(F_a)=0.2959(1.96)^2+2\times0.8520(1.2852)^2=3.9511,
\]
\[
 Q^F_{1,3}(F_a)=0.2959(1.96)^3+2\times0.8520(1.2852)^3=5.8449.
\]
Repitiendo para $F_b$ (intercambiando los roles de $x,y$, as\'i que en
efecto $d\cdot Cd$ recoge el $0.81$ y el $1.96$ en el patr\'on opuesto)
se obtiene
\[
 Q^F_{1,2}(F_b)=3.95107,\qquad Q^F_{1,3}(F_b)=5.7357.
\]
Ahora compara los dos estados directamente:
\[
 \big|Q^F_{1,2}(F_a)-Q^F_{1,2}(F_b)\big| \approx 0.00001,
 \qquad
 \big|Q^F_{1,3}(F_a)-Q^F_{1,3}(F_b)\big| \approx 0.109 .
\]
Este es el n\'umero m\'as importante de todo este documento.
\textbf{La feature de potencia 2 apenas distingue $F_a$ de $F_b$}
(una diferencia en el quinto decimal, para este sistema) --resulta
que un agregado \emph{cuadr\'atico} balanceado en 2D obedece una
simetr\'ia oculta casi exacta al intercambiar los dos ejes, as\'i que es
un detector de anisotropy casi ciego, no muy distinto de $I_1$ mismo en
este aspecto. \textbf{La feature de potencia 3, en cambio, difiere en
aproximadamente un $1.9\%$ de su propio valor} --una diferencia real,
estructural, no accidental. Esto es exactamente por qu\'e el modelo usa
t\'erminos tanto de $p=2$ como de $p=3$ para cada direction system, y
por qu\'e el paper principal dice que esto ``no es cosm\'etico'': sin
los t\'erminos c\'ubicos, el modelo estar\'ia mucho m\'as cerca de ser
isotropic de lo previsto, casi por accidente, por razones que no tienen
nada que ver con el entrenamiento y todo que ver con una coincidencia
geom\'etrica de 2D en potencias pares.
\end{worked}

¿Por qu\'e \emph{tres} direction systems separados, en lugar de uno?
Un solo sistema de tres direcciones ya rompe la isotropy (el Ejemplo
\ref{ex:worked-full} lo demuestra), pero solo muestrea la respuesta
direccional del material en tres \'angulos espec\'ificos. La rigidez de
un material anisotropic real puede variar de forma continua y
complicada con la direcci\'on; tres sistemas independientes, con
\'angulos distintos entre s\'i (nueve direcciones en total, cada
sistema individualmente balanceado), le dan al modelo muchas m\'as
``sondas'' independientes de rigidez direccional para combinar,
exactamente como usar m\'as sensores en m\'as \'angulos da una imagen
m\'as rica de una se\~nal direccional desconocida que un solo trio de
sensores.

\section{¿De d\'onde sali\'o esto -- lo inventamos nosotros?}
\label{sec:provenance}

Separ\'emoslo honestamente en dos capas, porque la respuesta honesta
tiene dos partes distintas.

\paragraph{La estrategia general (no es nuestra).} Construir una
energ\'ia \emph{polyconvex} como una input-convex neural network (ICNN)
aplicada a \emph{invariants} elegidos a mano, cada uno individualmente
\emph{polyconvex}, es el enfoque de Klein et
al.~\cite{Klein2022}, que a su vez se apoya en la teor\'ia original de
existencia de Ball~\cite{Ball1977} y en la tradici\'on cl\'asica de
hyperelasticity basada en \emph{invariants} que se remonta a Ogden y
antes~\cite{Ogden1984}. Klein et al.\ atan sus \emph{invariants} a un
\emph{named} material symmetry group (c\'ubico, transversalmente
isotropic, \dots) mediante un structural tensor cl\'asico --una t\'ecnica
bien establecida, de libro de texto (structural-tensor / integrity-basis
representation theory para \emph{invariants} anisotropic) aplicada por
muchos autores antes que ellos.

\paragraph{El mecanismo espec\'ifico (s\'i es nuestro).} Lo que es
genuinamente construcci\'on propia de este proyecto es la decisi\'on de
\emph{no} atar los \emph{invariants} a ning\'un named symmetry group en
absoluto: en cambio, usar varios direction systems gen\'ericos,
individualmente balanceados v\'ia~\eqref{eq:balance}, sin ninguna
clausura de rotaci\'on de $90^\circ$ ni ning\'un promedio $D_4$
incorporado en ninguna parte. Esta es la respuesta espec\'ifica al
problema abierto que Linden et al.~\cite{Linden2023} se\~nalan
expl\'icitamente en su propio paper --que puede no existir un conjunto
de \emph{invariants} completo y compatible con \emph{polyconvexity} para
una symmetry class general, desconocida-- y es lo que le permite a esta
construcci\'on reclamar \emph{objectivity} exacta, un \emph{stress-free
reference state} exacto, y \emph{polyconvexity} exacta \emph{sin nunca
asumir ni verificar} a qu\'e symmetry class, si acaso alguna, pertenece
el material. Los \'angulos y pesos num\'ericos espec\'ificos de la
tabla de direction systems son una soluci\'on v\'alida de la
\emph{balance condition}~\eqref{eq:balance} (que est\'a subdeterminada)
--hay infinitas soluciones v\'alidas, y las que se usaron no se derivaron
de ning\'un principio f\'isico m\'as profundo, m\'as all\'a de satisfacer
la ecuaci\'on~\eqref{eq:balance} con precisi\'on de m\'aquina y de dar
tres patrones angulares visiblemente distintos, no relacionados por
$D_4$; si quieres una elecci\'on m\'as limpia o mejor optimizada, la
\emph{balance condition} es la \'unica restricci\'on real que hay que
preservar.

\section{Un resumen corto, en un p\'arrafo}
\label{sec:recap}

La \emph{convexity} ordinaria en $F$ es imposible para cualquier
energ\'ia \emph{objective} (Secci\'on~\ref{sec:naive-fails}). La
soluci\'on de Ball, \emph{polyconvexity}, pide \emph{convexity} no en
$F$ sino en la tripleta ampliada y ``levantada'' $(F,\Cof F,\det F)$,
tratada como variables independientes
(Secci\'on~\ref{sec:lifting-trick}) --exactamente el mismo truco que
promover $\det F$ a su propia coordenada libre en el
Ejemplo~\ref{ex:lifting}. Los \emph{invariants} cl\'asicos $I_1,I_2,I_3$
son la forma de libro de texto de construir energ\'ias
\emph{polyconvex} a partir de esta tripleta
(Secci\'on~\ref{sec:classical-invariants}), pero son \emph{isotropic}
por construcci\'on y demostrablemente ciegos a la direcci\'on
(Ejemplo~\ref{ex:isotropy-blind}). Nuestras propias features son una
generalizaci\'on directa y literal de $I_1$ y del \emph{invariant}
basado en $\Cof$ --m\'as direcciones, pesos generales que satisfacen una
\emph{balance condition}, y, de forma crucial, potencias tanto
cuadr\'aticas \emph{como} c\'ubicas, porque la potencia cuadr\'atica
sola es un detector de anisotropy casi ciego en 2D
(Ejemplo~\ref{ex:worked-full}). La estrategia general de
ICNN-sobre-\emph{invariants} es de Klein et al.; la elecci\'on
espec\'ifica de usar direction systems gen\'ericos, libres de symmetry,
en lugar de un structural tensor con nombre, es la contribuci\'on
propia de este proyecto, construida para responder a un problema
abierto planteado expl\'icitamente en la literatura a la que responde.

\begin{thebibliography}{9}
\bibitem{Ball1977}
J.~M. Ball. Convexity conditions and existence theorems in nonlinear
elasticity. \emph{Archive for Rational Mechanics and Analysis}, 63:337--403,
1977.
\bibitem{Ciarlet1988}
P.~G. Ciarlet. \emph{Mathematical Elasticity, Volume I: Three-Dimensional
Elasticity}. North-Holland, 1988.
\bibitem{Ogden1984}
R.~W. Ogden. \emph{Non-Linear Elastic Deformations}. Ellis Horwood,
Chichester, 1984 (reprinted by Dover, 1997).
\bibitem{Klein2022}
D.~K. Klein, M. Fern\'andez, R.~J. Martin, P. Neff, and O. Weeger.
Polyconvex anisotropic hyperelasticity with neural networks.
\emph{Journal of the Mechanics and Physics of Solids}, 159:104703, 2022.
\bibitem{Linden2023}
L. Linden, D.~K. Klein, K.~A. Kalina, J. Brummund, O. Weeger, and
M. K\"astner. Neural networks meet hyperelasticity: a guide to enforcing
physics. \emph{Journal of the Mechanics and Physics of Solids},
179:105363, 2023.
\end{thebibliography}

\end{document}
